\section{The shape of the proof}\label{sec:overview}

The top of the proof is short. \Cref{st:rh} follows from
\cref{st:grh} for the trivial character by the functional equation, in a few dozen lines on
Mathlib's interface. \Cref{st:grh} is obtained for every character from a statement
\lean{ZeroFreeAll} about primitive characters, and that from a theorem with three hypotheses,
\begin{gather*}
\text{\lean{every_character_of_green_v_win}}\colon\\
\mathsf{Green}\ \longrightarrow\ \mathsf{VWin}\ \longrightarrow\ \mathsf{Nonempty}(\mathsf{HypsAlt})
\ \longrightarrow\ \mathsf{EveryCharacter},
\end{gather*}
each of which is in turn proved in the archive. They are the three inputs of the argument:
\begin{itemize}
\item $\mathsf{Green}$: a statement about the characters of finite general linear groups, in
  the parametrisation of J.~A.~Green \cite{green1955}, used to count local types;
\item $\mathsf{VWin}$: a bound on Galois cohomology groups attached to an auxiliary field,
  valid on windows of scales (\cref{sec:cohomology});
\item $\mathsf{HypsAlt}$: a record of forty-six fields about families of lattices of
  automorphic forms on unitary groups (\cref{sec:lattices,sec:frame,sec:literature}).
\end{itemize}

The body of that theorem \lvP{} applies a windowed form of Landau's theorem with window
exponent $\lambda=\tfrac32$ to a statement $\mathsf{TopG}$, which says: for every small
$\eps>0$ there is a family of arithmetic data, indexed by scales $x$ in long windows, on which
a condition \lean{ChainTied} holds. \lean{ChainTied} packages eleven inequalities between real
numbers (\cref{sec:lattices}) and implies $G(x)\le x^{1/2+C\eps}$ by an argument that involves
no arithmetic beyond an identity of character sums (\cref{sec:analytic}).

So the proof has two halves that meet in eleven real inequalities:
\begin{center}
\begin{tabular}{@{}p{0.46\textwidth}@{\qquad}p{0.46\textwidth}@{}}
\toprule
\emph{Analytic half} (\cref{sec:analytic}) & \emph{Arithmetic half} (Sections~\ref{sec:lattices}--\ref{sec:hosts})\\
\midrule
Landau on windows; an exact identity for a count of prime powers; averaging; balancing three
error terms at $m\approx\sqrt x$.
&
Lattices of forms, period maps, divisibility at inert primes, a slope inequality, a period
computation, a trace formula, a cohomological bound, auxiliary fields.\\
\bottomrule
\end{tabular}
\end{center}
The analytic half is elementary and a few hundred lines long. Almost all of the archive is the
arithmetic half and the published theorems it formalises.

\subsection*{Notation}
Throughout, $x$ is a large real scale and $\eps>0$ is small and fixed; $b=\exp(x^{\eps})$;
$m$ is the number of lifts, $x^{1/2-\eps}\le m\le x^{1/2+\eps/2}$; $E$ is the auxiliary
(``host'') field, with Galois group $\mathfrak G=(\Z/2)^2\times H$ over $\Q$ and $n=\abs H$;
$\Phi\subset\mathfrak G$ is a CM type; $\gamma$ is Euler's constant, $\gamma_p=p/(p^2-1)$ is
the rate of \cref{sec:adm,sec:divisibility}, and $\gamma_i$ are the weights of a frame
(\cref{sec:slope,sec:frame}). We write $\ph(y)=y\log y-y$, so
that $\ph(a)=\log a!+O(\log a)$.
